%! AP Statistics — Unit 9, Lesson 9.1 — Warm-Up
\documentclass[10pt]{article}
\usepackage{apstats-article}
\usepackage{apstats-boxes}

\begin{document}

\pageheader{Unit 9, Lesson 9.1}{Warm-Up}
\namedateperiod

\vspace{0.15in}

A researcher records the number of hours students studied ($x$) and their score
on a statistics exam ($y$) for a random sample of 30 students. The LSRL is:
\[
    \hat{y} = 48.3 + 4.2x
\]

\begin{enumerate}

  \item Identify the slope and $y$-intercept of the LSRL. Interpret the
        \textbf{slope} in context.

        Slope: \blank{2cm} \quad $y$-intercept: \blank{2cm}

        \vspace{0.05in}
        Interpretation: \blank{9cm}

        \vspace{0.1in}
        \blank{11cm}

        \vspace{0.1in}

  \item Use the LSRL to predict the exam score for a student who studied
        for 5 hours. Show your calculation.

        \vspace{0.05in}
        $\hat{y} = 48.3 + 4.2($\blank{1.5cm}$) = $ \blank{2.5cm}

        \vspace{0.1in}

  \item The slope $b = 4.2$ comes from one sample of 30 students. If we took a
        different random sample of 30 students from the same population, would
        we get exactly $b = 4.2$ again? What does this suggest about the
        relationship between the sample slope and the population slope $\beta$?

        \writelines{3}

\end{enumerate}

\end{document}
